% Классификация по manual/chapter-3-algorithms.tex, sec:planar.
\begin{tikzpicture}[course diagram,x=1mm,y=1mm,
  condition/.style={draw=courserule,fill=coursetint,line width=.6pt,
    text width=43mm,minimum height=10mm,align=center,inner sep=1.5mm,
    font=\normalsize},
  result/.style={draw=courserule,fill=coursetint,line width=.6pt,
    text width=76mm,minimum height=12mm,align=center,inner sep=1.5mm,
    font=\normalsize},
  branch/.style={draw=courseaccent,line width=.7pt,->}]
  \path[use as bounding box] (0,0) rectangle (148,42);
  \node[condition] (tangent-test) at (26,35) {$|\Delta a|=\Delta e>0$};
  \node[condition] (disjoint-test) at (26,21) {$|\Delta a|>\Delta e$};
  \node[condition] (crossing-test) at (26,7) {$|\Delta a|<\Delta e$};
  \node[result] (tangent) at (108,35)
    {Касающиеся орбиты\\
     {\fontsize{9.5}{11.5}\selectfont Один трансверсальный импульс}};
  \node[result] (disjoint) at (108,21)
    {Непересекающиеся орбиты\\
     {\fontsize{9.5}{11.5}\selectfont Два импульса одного знака}};
  \node[result] (crossing) at (108,7)
    {Пересекающиеся орбиты\\
     {\fontsize{9.5}{11.5}\selectfont Два импульса противоположных знаков}};
  \draw[branch] (tangent-test.east)--(tangent.west);
  \draw[branch] (disjoint-test.east)--(disjoint.west);
  \draw[branch] (crossing-test.east)--(crossing.west);
\end{tikzpicture}
